


Generated by \texttt{python -m crashguard.frame\_layout\_doc} from \texttt{src/crashguard/frame\_layout.py} (the field tables) and the stored keys. \texttt{tests/test\_frame\_layout.py} checks the CODE rows against \texttt{packet.py} and \texttt{sct\_protocol.py} and the DESIGN rows against the reference functions. This document closes blocker B5 of the final verification run (the byte layout of the freshness field and the boot counter).

\textbf{Status of each row.} CODE: an encoder in the repository fixes it (\texttt{packet.py}, \texttt{sct\_protocol.py}, \texttt{firmware\_rd/cg\_rd.c}, \texttt{firmware\_r6/cg\_r6.c}). DESIGN: chosen in the delta verification run to close the gap; no encoder in the repository writes it, no stored result depends on it, and no frame size of an evaluated design changes. A choice of the author overrides a DESIGN row. The firmware does not implement any of this (no physical hardware was used).

\textbf{Conventions.} Little-endian for every multi-byte field of every frame. The tag is the leading bytes of HMAC-SHA256 under the per-vehicle key K\_vehicle = HMAC(master, vehicle id) (R4), computed over every byte that precedes it, and compared without early exit. The micro-confirm tag input uses a big-endian boot counter (\texttt{cg\_micro\_tag}); it is a MAC input, not a frame field.

\subsection*{1. The five frames of the two-stage alert}

{\scriptsize\raggedright
\begin{longtable}{p{0.237\linewidth}p{0.237\linewidth}p{0.237\linewidth}p{0.237\linewidth}}
\toprule
Frame & Bytes & Tag bytes & Sent when \\
\midrule
\endhead
Provisional flash & 14 (evaluated) or 15 (with the boot byte, recommended) & 8 & at the early-exit decision \\
Full alert & 97 & 8 & at the full-window decision, or alone when there was no early alarm \\
ACK & 26 & 8 & by the base, for a flash (acked event id 0xFFFFFFFF) or an alert \\
Cancel & 14 or 15 & 8 & at the full-window decision when the full window does not alarm; three transmissions 0.3 s apart (ns-3) without acknowledgement \\
Micro-confirm and micro-cancel & 8 (design); 13 in ns-3 & 5 & after the flash was acknowledged and while the base holds its pending entry (30 s) \\
\bottomrule
\end{longtable}}

\subsubsection*{2.1 Full alert, 97 bytes}

{\scriptsize\raggedright
\begin{longtable}{p{0.158\linewidth}p{0.158\linewidth}p{0.158\linewidth}p{0.158\linewidth}p{0.158\linewidth}p{0.158\linewidth}}
\toprule
Offset & Size & Field & In the tag & Source & Meaning \\
\midrule
\endhead
0 & 1 & version & yes & CODE & packet.VERSION = 1 \\
1 & 1 & msg\_type & yes & CODE & 1 = alert (domain separation: the type is inside the tagged bytes) \\
2 & 4 & node\_id & yes & CODE & u32 \\
6 & 4 & event\_id & yes & CODE & u32, unique per event of the vehicle \\
10 & 4 & freshness field & yes & DESIGN & was \texttt{counter} (u32). Time-stamped form: GPS seconds since 1980-01-06 00:00:00 UTC, bit 31 = 0. No-time form: bit 31 = 1, bits 30 to 16 = epoch (boot counter mod 2\textasciicircum{}15), bits 15 to 0 = counter \\
14 & 4 & timestamp\_ms & yes & CODE & u32, unchanged \\
18 & 1 & severity & yes & CODE & 0 low to 3 critical \\
19 & 1 & impact\_side & yes & CODE & 0 F, 1 R, 2 L, 3 Rt, 4 V \\
20 & 2 & angle\_deg x10 & yes & CODE & u16 \\
22 & 4 & lat x1e7 & yes & CODE & i32 \\
26 & 4 & lon x1e7 & yes & CODE & i32 \\
30 & 2 & peak\_g x100 & yes & CODE & u16 \\
32 & 2 & ml\_prob x10000 & yes & CODE & u16 \\
34 & 1 & rule\_sev & yes & CODE & u8 \\
35 & 1 & sw\_mask & yes & CODE & u8 \\
36 & 1 & sw\_first & yes & CODE & u8 \\
37 & 1 & rollover & yes & CODE & u8 \\
38 & 2 & delta\_v\_kmh x100 & yes & CODE & u16 \\
40 & 2 & asi x100 & yes & CODE & u16 \\
42 & 1 & sequence in second & yes & DESIGN & u8, the lead byte of the 47 reserved bytes: number of alerts the vehicle already sent with the same freshness field (0 for the initial alert; saturates at 255) \\
43 & 4 & boot counter & yes & DESIGN & u32, non-volatile, increased and written once per boot (four of the reserved bytes) \\
47 & 42 & reserved & yes & CODE & zeros; the 47 reserved bytes of packet.RESERVED\_LEN minus the 5 bytes used above \\
89 & 8 & tag & is the tag & CODE & HMAC-SHA256(K\_vehicle, bytes 0 to 88)[:8] \\
\bottomrule
\end{longtable}}

\subsubsection*{2.2 ACK, 26 bytes}

{\scriptsize\raggedright
\begin{longtable}{p{0.158\linewidth}p{0.158\linewidth}p{0.158\linewidth}p{0.158\linewidth}p{0.158\linewidth}p{0.158\linewidth}}
\toprule
Offset & Size & Field & In the tag & Source & Meaning \\
\midrule
\endhead
0 & 1 & version & yes & CODE & 1 \\
1 & 1 & msg\_type & yes & CODE & 2 = ACK \\
2 & 4 & node\_id & yes & CODE & u32, the vehicle that is acknowledged (the ACK is tagged with that vehicle's key) \\
6 & 4 & acked event\_id & yes & CODE & u32; 0xFFFFFFFF = the ACK of a provisional frame (sct\_protocol.FLASH\_EVENT\_ID) \\
10 & 4 & acked freshness field & yes & DESIGN & was \texttt{acked counter}: the 32-bit freshness field of the acknowledged frame (for a flash: reconstructed by the base) \\
14 & 4 & base counter & yes & CODE+DESIGN & u32, the base's own counter; no model uses it. In the ACK of a provisional frame its low byte is the pending index that the micro frame returns (DESIGN, finalize run) \\
18 & 8 & tag & is the tag & CODE & HMAC-SHA256(K\_vehicle, bytes 0 to 17)[:8] \\
\bottomrule
\end{longtable}}

\subsubsection*{2.3 Provisional flash, 14 bytes (the evaluated frame)}

{\scriptsize\raggedright
\begin{longtable}{p{0.158\linewidth}p{0.158\linewidth}p{0.158\linewidth}p{0.158\linewidth}p{0.158\linewidth}p{0.158\linewidth}}
\toprule
Offset & Size & Field & In the tag & Source & Meaning \\
\midrule
\endhead
0 & 1 & version and type & yes & CODE & high nibble version (1), low nibble type: 4 = flash (domain separation inside the tagged bytes) \\
1 & 2 & short node id & yes & CODE & u16, low 16 bits of node\_id; ASSUMPTION: unique in the fleet \\
3 & 2 & freshness field, low 16 bits & yes & CODE+DESIGN & CODE: low 16 bits of the 32-bit counter of the event. DESIGN: of the 32-bit freshness field of the event (the flash and the full alert of one event carry the same field) \\
5 & 1 & flags and event type & yes & CODE+DESIGN & bit 7 (DESIGN) = no-time form; bits 6 to 4 = 0: the provisional frame carries no severity class (the encoder allows a class there; the design leaves it zero); bits 3 to 0 = event type (0 crash, 1 rollover, 2 other) \\
6 & 8 & tag & is the tag & CODE & HMAC-SHA256(K\_vehicle, bytes 0 to 5)[:8] \\
\bottomrule
\end{longtable}}

\subsubsection*{2.4 Cancel, 14 bytes}

{\scriptsize\raggedright
\begin{longtable}{p{0.158\linewidth}p{0.158\linewidth}p{0.158\linewidth}p{0.158\linewidth}p{0.158\linewidth}p{0.158\linewidth}}
\toprule
Offset & Size & Field & In the tag & Source & Meaning \\
\midrule
\endhead
0 & 1 & version and type & yes & CODE & type 5 = cancel \\
1 & 2 & short node id & yes & CODE & as the flash \\
3 & 2 & freshness field, low 16 bits & yes & CODE+DESIGN & of the event that is cancelled \\
5 & 1 & flags and reason & yes & CODE+DESIGN & bit 7 = no-time form (DESIGN); bits 3 to 0 = reason (0 = driver cancel) \\
6 & 8 & tag & is the tag & CODE & HMAC-SHA256(K\_vehicle, bytes 0 to 5)[:8] \\
\bottomrule
\end{longtable}}

\subsubsection*{2.5 Provisional flash with the boot byte, 15 bytes (recommended)}

{\scriptsize\raggedright
\begin{longtable}{p{0.158\linewidth}p{0.158\linewidth}p{0.158\linewidth}p{0.158\linewidth}p{0.158\linewidth}p{0.158\linewidth}}
\toprule
Offset & Size & Field & In the tag & Source & Meaning \\
\midrule
\endhead
0 & 1 & version and type & yes & CODE & as FLASH-14 \\
1 & 2 & short node id & yes & CODE & as FLASH-14 \\
3 & 2 & freshness field, low 16 bits & yes & CODE+DESIGN & as FLASH-14 \\
5 & 1 & flags and event type & yes & CODE+DESIGN & as FLASH-14 \\
6 & 1 & boot counter, low 8 bits & yes & DESIGN & u8: B mod 256; the base rebuilds the 32-bit value nearest above its last known one \\
7 & 8 & tag & is the tag & CODE & HMAC-SHA256(K\_vehicle, bytes 0 to 6)[:8] \\
\bottomrule
\end{longtable}}

\subsubsection*{2.6 Cancel with the boot byte, 15 bytes}

{\scriptsize\raggedright
\begin{longtable}{p{0.158\linewidth}p{0.158\linewidth}p{0.158\linewidth}p{0.158\linewidth}p{0.158\linewidth}p{0.158\linewidth}}
\toprule
Offset & Size & Field & In the tag & Source & Meaning \\
\midrule
\endhead
0 & 1 & version and type & yes & CODE & type 5 \\
1 & 2 & short node id & yes & CODE & as FLASH-14 \\
3 & 2 & freshness field, low 16 bits & yes & CODE+DESIGN & of the cancelled event \\
5 & 1 & flags and reason & yes & CODE+DESIGN & as CANCEL-14 \\
6 & 1 & boot counter, low 8 bits & yes & DESIGN & as FLASH-15 \\
7 & 8 & tag & is the tag & CODE & HMAC-SHA256(K\_vehicle, bytes 0 to 6)[:8] \\
\bottomrule
\end{longtable}}

\subsubsection*{2.7 Micro-confirm and micro-cancel, 8 bytes}

{\scriptsize\raggedright
\begin{longtable}{p{0.158\linewidth}p{0.158\linewidth}p{0.158\linewidth}p{0.158\linewidth}p{0.158\linewidth}p{0.158\linewidth}}
\toprule
Offset & Size & Field & In the tag & Source & Meaning \\
\midrule
\endhead
0 & 1 & pending index & is the tag & DESIGN & u8, assigned by the base when it acknowledges the provisional frame (low byte of the \texttt{base counter} of that ACK) and sent back here; it names the pending entry (vehicle key, T1, boot counter), so the base verifies the tag under one candidate. It is not in the MAC input: a changed index selects another entry, whose tag then fails (probability 2\textasciicircum{}-40 per frame to pass). F3 of the finalize run \\
1 & 1 & type, class, event, flag & yes & CODE+DESIGN & bit 7 = type (0 confirm, 1 cancel); bits 6 to 4 = severity class (for a cancel the reason); bits 3 to 2 = event type; bit 1 = position-offset flag; bit 0 = 0 (spare). V17 section 3 fixes the fields, the bit positions are DESIGN. The tag function treats the byte as opaque (cg\_micro\_tag) \\
2 & 1 & position offset & yes & CODE+DESIGN & high nibble east, low nibble north, each signed 4 bits in units of 128 m relative to the last position the base holds for the vehicle; 0x00 when the flag is 0 \\
3 & 5 & tag & is the tag & CODE & HMAC-SHA256(K\_vehicle, 'M' followed by T1 (8 bytes), the two content bytes (the type byte and the position offset) and the boot counter (u32, big endian))[:5]; T1 = the 8-byte tag of the provisional frame; the frame carries neither T1 nor the vehicle, the index names them. Bytes 1 to 7 are the 7-byte frame of \texttt{cg\_micro\_pack} \\
\bottomrule
\end{longtable}}

\subsubsection*{2.8 The freshness field}

\begin{itemize}
\item \textbf{Time-stamped form} (the vehicle has a GPS fix): bit 31 = 0; bits 30 to 0 = seconds since 1980-01-06 00:00:00 UTC as the receiver reports them (the base uses the same scale; the window W = 10 s is not affected by the leap-second offset if both sides use one scale). Valid until the year 2116.
\item \textbf{No-time form} (no fix): bit 31 = 1; bits 30 to 16 = epoch = boot counter modulo 2\textasciicircum{}15, written once per boot; bits 15 to 0 = counter of alerts since boot.
\item \textbf{Identifier of an alert} for the freshness rule: the triple (sender, field, sequence in second); with the boot counter, (sender, boot, field, sequence). The base accepts a time-stamped alert if |field - now| <= W and the triple is not in its ring of R entries; this is \texttt{cg\_fresh\_accept(ring, field, seq, now, W, R)}.
\item \textbf{In the flash and the cancel} only the low 16 bits are sent, with the no-time flag. For a time-stamped frame the base rebuilds the 32-bit value that is congruent to those 16 bits modulo 2\textasciicircum{}16 and nearest its own clock (\texttt{reconstruct\_field}; unambiguous while the sender's clock is within 2\textasciicircum{}15 s of the base's, and W is 10 s). For a no-time frame it uses the epoch it holds for the vehicle.
\item \textbf{Two flashes of one vehicle inside one GPS second} carry the same 16 bits and no sequence: the second is treated as a repeat of the earlier one and is acknowledged again, not processed again. The full alert carries the sequence, so no alert is lost.
\item \textbf{Boot counter.} 32 bits in non-volatile memory, increased and written once per boot (the write the boot-epoch design already counts). The alert carries all 32 bits; the 15-byte flash and cancel carry the low 8 bits, and the base rebuilds the 32-bit value as the smallest value not below the last one it holds for the vehicle (\texttt{reconstruct\_boot}). The 14-byte frames of the evaluated design carry no boot byte; they are safe only if a reboot takes at least one second (the ProVerif finding, \texttt{FORMAL\_RERUN.md}).
\end{itemize}

\subsection*{3. Tag lengths and domain separation}

{\scriptsize\raggedright
\begin{longtable}{p{0.237\linewidth}p{0.237\linewidth}p{0.237\linewidth}p{0.237\linewidth}}
\toprule
Frame & Tag & Separation inside the tagged bytes & Letter in the formal models \\
\midrule
\endhead
alert & 8 bytes & msg\_type byte = 1 at offset 1, inside the tagged bytes; length 97 & 'A' \\
ACK & 8 bytes & msg\_type byte = 2 at offset 1, inside the tagged bytes; length 26 & 'ACK' \\
provisional frame & 8 bytes & type nibble 4 in byte 0, inside the tagged bytes; length 14 or 15 & 'F' (ProVerif and Tamarin R6), 'P' (round-4 model) \\
cancel (14 or 15 bytes) & 8 bytes & type nibble 5 in byte 0, inside the tagged bytes & no letter in the frozen models; micro-cancel 'X' \\
micro-confirm and micro-cancel & 5 bytes & ASCII 'M' (0x4D) in front of the MAC input; confirm and cancel differ by bit 7 of byte 0, inside the tagged bytes; length 8 & 'M' (confirm), 'X' (cancel) \\
\bottomrule
\end{longtable}}

The letters of the models are abstractions of the type values and lengths above (row F15 of section 8). Forgery effort per attempt: 2\textasciicircum{}-64 for an 8-byte tag and 2\textasciicircum{}-40 for the 5-byte micro tag, which the pending index keeps at one candidate entry (ANALYTICAL; section 7).

\subsection*{4. Pending table and the lookup of a micro frame}

The base keeps a pending entry (vehicle key, T1, boot counter) for every provisional frame it acknowledged, for 30 s (the frozen lifetime of R6-1, \texttt{r6.r1.ns3.params.pend\_s}), in a table slot of one byte. The ACK of the provisional frame returns the slot number in the low byte of its \texttt{base counter}; the micro frame carries it as its leading byte. The base verifies the micro tag under the one entry that the index names (\texttt{micro\_lookup}), and treats the frame as the confirm or cancel of that entry. A frame whose index is unknown or whose tag fails (an expired entry, a car that rebooted so that the base's boot counter is old, or a changed index) is rejected, and the node's retry then sends the full frame (the fallback of V17 section 3). The 7-byte lookup by trial (\texttt{micro\_lookup\_trial}) was rejected: it costs one tag computation per pending entry and multiplies the forgery probability by their number. The micro frame is a smaller frame, not a different authority: the full alert still carries the content.

\subsection*{5. The 8-byte design against the 13-byte ns-3 frame}

The LoRaWAN module of ns-3 sends at least 13 bytes (its header with the 4-byte device address and the message integrity code), so the simulated micro frame has 13 bytes (\texttt{R6\_MICRO\_BYTES}). The design frame has 8 bytes (the 7-byte frame of \texttt{cg\_micro\_pack} behind the index), so it stays within the evaluated size and the airtime results are conservative for it. The airtime in section 7 shows that at SF7 the 13-byte frame costs the airtime of the 14-byte flash, so against the 14-byte flash the micro frame saves nothing in the simulation; the saving is against the 97-byte confirm. The design frame would be shorter: 36.1 ms against 46.3 ms at SF7.

\subsection*{6. Working flow of one event, and what happens to an unacknowledged flash}

\begin{itemize}
\item Trigger at t = 0. The early-exit rule decides at t\_E (median 50 ms, mean 112 ms on the held-out events) and, if it alarms, the node queues the flash (ready at t\_E + 8 ms).
\item The full window (1.8 s plus 8 ms) decides at T\_F = 1.808 s: a confirm (97-byte alert, or the micro-confirm if the flash was acknowledged and the entry is pending) or three cancels (14 bytes, 0.3 s apart, unacknowledged).
\item The base accepts a flash (state PRE: crash reported, severity pending), a confirm (CONFIRMED, with or without a flash) or a cancel (CANCELLED, only if the base holds that event).
\end{itemize}

**What the code does with an unacknowledged flash (read from \texttt{xsim\_alert.cc} and \texttt{r3\_t1.py}).\textbf{ In the evaluated two-stage policies (r3t, r6m) and in the trace replays the node keeps one FIFO queue per event: [flash, confirm] or [flash, cancel, cancel, cancel]. The flash stays at the head of the queue until its acknowledgement arrives or its 4 tries are used (each try is followed by the two receive windows, 0.32 s in ns-3). Nothing is withdrawn when the full decision is ready. }For a cancel:\textbf{ the flash is still retried to the end of its stage and the cancel is sent after it, whether or not the flash was ever delivered (the base rejects a cancel for an event it does not hold). }For a confirm:\textbf{ the confirm does not replace the flash; both are sent, the flash stage ahead of the confirm, and in the micro policy r6m the node sends the full 97-byte confirm when the flash was not acknowledged. }Only the policy r3a** (24-byte provisional frame, V15, not supported) preempts: a ready confirm or cancel takes the radio ahead of the provisional repeats, a due cancel withdraws the queued provisional frames, and the acknowledgement of the confirm ends the provisional repetition.

\textbf{Design recommendation (not implemented in r3t or r6m): preempt.} When the full decision is ready, withdraw the remaining flash tries and send the decision's frames. 
ANALYTICAL saving, from the time-on-air formula: the flash tries start at s\_f + j (A14 + 0.32 s), j = 0 to 3; the tries that start at or after T\_F are the ones that preemption removes.

{\scriptsize\raggedright
\begin{longtable}{p{0.190\linewidth}p{0.190\linewidth}p{0.190\linewidth}p{0.190\linewidth}p{0.190\linewidth}}
\toprule
SF & A14 (ms) & tries that start at or after T\_F (of 4) & airtime saved per event with an unacknowledged flash (s) & end of the 4th try (s) \\
\midrule
\endhead
7 & 46.34 & 0 & 0.000 & 1.203 \\
8 & 82.43 & 0 & 0.000 & 1.348 \\
9 & 164.86 & 0 & 0.000 & 1.677 \\
10 & 288.77 & 1 & 0.289 & 2.173 \\
11 & 659.46 & 2 & 1.319 & 3.656 \\
12 & 1155.07 & 2 & 2.310 & 5.638 \\
\bottomrule
\end{longtable}}

At SF7 to SF9 all four tries end before T\_F (the 4th ends at 1.68 s at SF9), so preemption saves nothing there and does not change the stored results. From SF10 it removes 1 try (the 4th would end at 2.17 s, after T\_F), and at SF11 and SF12 2 tries, which also delays the confirm behind the flash stage.

**Formal check (FORMAL, bounded; \texttt{python -m crashguard.vf2\_preempt}, keys \texttt{vf2.preempt.*}).\textbf{ An untimed model with one event, one base, every try lost, delivered with a lost ACK, or delivered and acknowledged, 1 to 4 flash tries before the decision, a confirm of up to 4 tries or three cancels, enumerated exhaustively for the FIFO node and the preempting node. Result: 8 (decision, tries) cases; the set of base states reachable at the end is the same for both nodes (preemption adds no terminal state); a delivered final-stage frame always takes the base out of PRE; the base can stay in PRE with no further frame }only if every frame of the final stage was lost** (all four confirm tries, or all three unacknowledged cancels), and this holds for the FIFO node as much as for the preempting one. Preemption therefore cannot leave the base in pre-alert forever by itself; the stuck state is a property of the unconfirmed cancel and of a lost confirm, and a base-side timeout (PRE to EXPIRED when neither confirm nor cancel arrives) removes it in both variants (\texttt{pre\_alert\_forever\_with\_timeout} = 0 for the cancel and 0 for the confirm). The code has no such timeout; it is a design recommendation. The model does not cover time, duty limits or several vehicles.

\subsection*{7. Airtime and forgery numbers (ANALYTICAL)}

{\scriptsize\raggedright
\begin{longtable}{p{0.119\linewidth}p{0.119\linewidth}p{0.119\linewidth}p{0.119\linewidth}p{0.119\linewidth}p{0.119\linewidth}p{0.119\linewidth}p{0.119\linewidth}}
\toprule
SF & 7 bytes & 8 bytes & 13 bytes & 14 bytes & 15 bytes & 26 bytes & 97 bytes \\
\midrule
\endhead
7 & 36.10 & 36.10 & 46.34 & 46.34 & 46.34 & 61.70 & 169.22 \\
8 & 72.19 & 72.19 & 82.43 & 82.43 & 92.67 & 113.15 & 297.47 \\
9 & 123.90 & 123.90 & 164.86 & 164.86 & 164.86 & 205.82 & 533.50 \\
10 & 247.81 & 247.81 & 288.77 & 288.77 & 329.73 & 411.65 & 985.09 \\
11 & 495.62 & 495.62 & 577.54 & 659.46 & 659.46 & 823.30 & 2134.02 \\
12 & 991.23 & 991.23 & 1155.07 & 1155.07 & 1155.07 & 1646.59 & 3940.35 \\
\bottomrule
\end{longtable}}

Times in ms (125 kHz, coding rate 4/5, explicit header, CRC, 8 preamble symbols). 13, 14 and 15 bytes have the same airtime at SF7, SF9, SF12; the 15-byte frame costs more airtime at the other spreading factors. The 13-byte ns-3 micro frame is 1.28 times the 8-byte design at SF7 and 1.33 times at SF9.

\textbf{Forgery with the pending index.} The base verifies one micro frame under one entry, so the chance that a random 5-byte tag passes is 2\textasciicircum{}-40 per frame, the single-vehicle value. At the attacker's ceiling of 36 s / A8 frames per hour at SF7 (the 1\% allowance), the expected years to one accepted forgery are 2\textasciicircum{}40 / frames per year (the formula of \texttt{r6.b.forge.tag5.expected\_years}, with the 8-byte frame time). The rejected lookup by trial, for comparison, tests n entries per frame:

{\scriptsize\raggedright
\begin{longtable}{p{0.317\linewidth}p{0.317\linewidth}p{0.317\linewidth}}
\toprule
Lookup and pending entries n & Probability per frame & Expected years to one forgery \\
\midrule
\endhead
pending index (the design): 1 entry & 9.09e-13 & 125,764 \\
trial lookup (rejected): 10 entries & 9.09e-12 & 12,576 \\
trial lookup (rejected): 60 entries & 5.46e-11 & 2,096 \\
trial lookup (rejected): 1000 entries & 9.09e-10 & 126 \\
\bottomrule
\end{longtable}}

The stored value \texttt{r6.b.forge.tag5.expected\_years} (125,764 years) used the 7-byte frame time; the 8-byte frame has the same time on air as the 7-byte frame at every spreading factor from 7 to 12 (36.10 ms at SF7), so the figure of the design is 125,764 years (the stored key is not changed). The pending index restores the single-vehicle value whatever the number of pending entries.

\subsection*{8. Reconciliation of the layout, the formal models and the ns-3 frames}

Rule for the paper: the proofs cover the protocol at the level of the stated message structure; the abstractions in this table are the ones that apply. Rows F1 to F5 are the five model-to-frame mismatches listed in \texttt{docs/04\_audits/FORMAL\_FIELD\_CHECK.md}, F6 is its gap; F7 to F16 are added here. No stored result was changed: every row is resolved by a change in a document or by a stated abstraction.

{\scriptsize\raggedright
\begin{longtable}{p{0.158\linewidth}p{0.158\linewidth}p{0.158\linewidth}p{0.158\linewidth}p{0.158\linewidth}p{0.158\linewidth}}
\toprule
ID & In the model or the simulation & In the layout & Kind & Resolution & What it limits \\
\midrule
\endhead
F1 & Vehicle = one name \texttt{v} in every model & 32-bit node id in the alert and the ACK; 16-bit short id in the flash and the cancel; a one-byte pending index and no vehicle id in the 8-byte micro frame & assumption and abstraction & State: uniqueness of the 16-bit short id inside the fleet is an assumption of the layout and is in no model; the model treats the id as an opaque unique value & Cross-vehicle statements hold only if short ids are unique; a collision is not modelled \\
F2 & Freshness id = one fresh name, or the pair (epoch, sequence), or the triple (boot, epoch, sequence) & Alert: 32-bit field at offset 10, sequence at offset 42, boot counter at 43 to 46. Flash and cancel: low 16 bits of the field plus a no-time flag. Layout was missing (blocker B5) & document changed & \texttt{FRAME\_LAYOUT.md} specifies the fields; the models treat each tuple as one value and do not model byte order or width & The lemmas say nothing about the reconstruction of the 32-bit field from 16 bits in the flash \\
F3 & Freshness window of two epochs (one or two GPS seconds) & Window W = 10 s in the simulations and the C ring (\texttt{cg\_fresh\_accept}) & abstraction & State: the lemmas (at most once per window) do not use the width, and the statement 'inside the window' is for the model's window & The 10 s window is not covered by a proof \\
F4 & Seen-set without bound & Ring of the last R entries (R up to 512), cleared at a base reboot & abstraction & State: a replay of an entry that the ring has dropped is outside the model; \texttt{ar.r1} measures it by simulation (H-R1b not supported at the held-out burst) & At-most-once is proved for an unbounded set only \\
F5 & Boot counter: a fresh value per boot, unbounded & 32-bit counter in non-volatile memory; the alert carries 4 bytes, the flash and the cancel carry the low 8 bits (15-byte frame), the micro tag input carries the full 32 bits (big endian) & document changed and abstraction & Layout adds the byte (15 bytes, airtime unchanged at SF7 and SF9); the model's unbounded value stands for the counter; the 8-bit wrap needs fewer than 256 boots between two frames the base hears & A car that reboots 256 times or more between two heard frames is outside the rule; not modelled \\
F6 & The freshness lemmas (Tamarin \texttt{crash\_alert\_*}) model the 97-byte alert and the 26-byte ACK only & Flash and cancel are 14 or 15 bytes & gap, narrowed & State: the round-6 micro-confirm models (ProVerif and Tamarin) add the flash tag, the micro-confirm and the micro-cancel; no model covers the freshness field of the flash or the cancel & No formal result for the freshness of the provisional frame \\
F7 & Micro-confirm in the model: \texttt{<'MICRO', v, T1, x, m>} carries the vehicle and T1 in clear; the base looks up the pending entry by (v, T1) & Micro frame (8 bytes) carries a one-byte pending index that the base assigned in the ACK of the provisional frame; the index names the entry (vehicle key, T1, boot counter) and the base verifies the tag under that one entry. ns-3: the 13-byte frame includes the LoRaWAN device address, so ns-3 looks up by address & abstraction & State: the model's (v, T1) stands for the entry that the index names; the index is outside the MAC input and chosen by the attacker in the model's worst case (the attacker picks (v, T1) already), so it gives the attacker nothing more. The proofs are for lookup by (v, T1); the index lookup has one candidate, so the forgery probability per frame is the single-entry value 2\textasciicircum{}-40 (section 7). The index variants of the ProVerif and Tamarin models (\texttt{formal/proverif/vf3\_micro\_confirm\_index.pv}, \texttt{formal/vf3\_micro\_confirm\_index.spthy}, outputs in \texttt{formal/out/vf3/}) carry a pending table keyed by the index and an attacker-chosen index: acceptance implies the vehicle sent that micro-confirm, binding to its provisional frame and the micro-cancel property are proved in both tools (\texttt{vf3.formal.*}) & The index lookup itself (a table slot, its reuse after expiry) is not modelled \\
F8 & Micro-cancel separated by the letter 'X' in the MAC input & Micro-confirm and micro-cancel share the letter 'M'; they differ in bit 7 of byte 0, which is inside the tagged bytes & abstraction & State: separation is by a bit that the tag covers, which has the same effect as a letter; the exact encoding is not modelled & A cross-use of confirm and cancel is excluded by the model, and by the layout through the type bit \\
F9 & Perfect MAC, no truncation & HMAC-SHA256 truncated to 8 bytes (5 bytes for the micro frame) & abstraction & State: forgery bounds 2\textasciicircum{}-64 and 2\textasciicircum{}-40 per attempt are ANALYTICAL and separate & No computational result \\
F10 & Not in a formal model & ns-3 micro frame: 13 bytes (the smallest LoRaWAN frame) against the 8-byte design & simulation against design & State: the 8-byte design is within the evaluated size; at SF7 and SF9 the 13-byte frame has the airtime of the 14-byte flash, so the ns-3 airtime of the micro scheme is higher than the design's (section 7), and the airtime results are conservative for the 8-byte frame; the saving is against the 97-byte confirm only & ns-3 results are conservative for the micro airtime \\
F11 & Hash ratchet K(i+1) = H(K(i), label), at most 3 steps (ProVerif, Tamarin) & No frame has a ratchet step field; the keys of the frames are the per-vehicle keys of R4 and the micro tag uses the boot counter, not a ratchet & gap & State: forward secrecy is a statement about the symbolic ratchet within three steps; an implementation needs a step index (derived from the boot counter and the count of acknowledged alerts, or a new field) & No claim of forward secrecy for the frames as laid out \\
F12 & Quarantine: counters per claimed id, alert kinds ok, implausible, repeat; 3 ids, 6 slots, 4 attacker frames per slot & The claimed id is in clear in every frame except the micro frame; the base code verifies the tag before anything else and has no quarantine counter & abstraction & State: the result is for the modelled mechanism (rule A is attackable, rule B holds by construction); counting after verification is a design rule, not an implemented check & No general result beyond the modelled counters \\
F13 & Per-vehicle key = perfect function kdf(master, v) & HMAC-SHA256(master, id), cached per new id & abstraction & State: symbolic function stands for the HMAC & No computational result \\
F14 & ACK = <'ACK', v, f, tag>, f one value & ACK names (node id, acked event id, acked freshness field) and carries a base counter & abstraction & State: f stands for the pair (event id, freshness field); the base counter is not modelled and enters no lemma & None for the lemmas stated \\
F15 & Domain separation by letters ('A', 'ACK', 'F', 'M', 'X', 'P') & By the type value inside the tagged bytes (alert 1, ACK 2, flash 4, cancel 5 in the nibble) and by lengths; micro: ASCII 'M' and the type bit & abstraction & State: each letter stands for the type value of section 3; the round-4 control without a letter (\texttt{crash\_alert\_prov\_nodomain}) is falsified as expected & Separation across frame kinds is proved for the letters, not for the byte encodings \\
F16 & No time in any model & Confirm due at 1.808 s, retries after the receive windows, hourly allowance & abstraction & State: models are untimed; timing results are from simulation and enumeration (\texttt{timed} models of round 3) & No proof about delivery before a deadline \\
F17 & Round-7 models: the alert carries the pseudonym p, the encrypted content and the tag; boot and counter are not sent & Layout rows of section 9 (DESIGN) & design and abstraction & State: the window of the next 16 pseudonyms is a table entry for every alert a vehicle sends; a vehicle outside the window or after a reboot is outside the proofs; the tag covers boot and counter as unsent associated data (the symbolic models carry them in the tag input or in the encryption nonce) & Unlinkability and the lemmas hold for resolved vehicles only \\
F18 & Pseudonym = pseudorandom function symbol, truncation not modelled & p = the leading L bytes of HMAC-SHA256; L = 2 (flash, cancel) or 4 (alert, ACK) & abstraction & State: truncation creates collisions with probability 0.216 per frame at L = 2 and 3.7e-06 at L = 4 (1,000 vehicles, m = 16); the tag decides a collision; collisions are not in the symbolic models & No claim of collision freedom at these lengths \\
F19 & Attacker holds no vehicle key, no timing or radio-level view & Position is encrypted in the alert; arrival time, frame length and radio fingerprint remain visible & out of scope & State: unlinkability is against a network attacker that reads, replays and injects frames; it is not claimed against timing correlation, length or fingerprint analysis & No claim beyond the symbolic attacker \\
F20 & Reserve policy of the node: enumeration model with integer costs and a bounded number of events & Retry caps of 4 tries, airtime of the 14-byte flash and 97-byte confirm & abstraction & State: Invariants 1 and 2 are checked exhaustively in the bounds of \texttt{ROUND7\_A.md} section 4, not proved for all sizes; Invariant 2 is conditional on the reserve being intact & No proof for unbounded numbers of events \\
\bottomrule
\end{longtable}}

\subsection*{9. Round 7 design: pseudonymous vehicle id (DESIGN)}

Round 7 (\texttt{docs/02\_preregistration/PREREGISTRATION\_V18.md}, \texttt{docs/03\_results/ROUND7\_B.md}). No firmware writes these fields; no stored result of an earlier round depends on them; the sizes of the frames do not change. The layout of sections 1 to 8 stays the layout of the evaluated design; this section is the proposed change.

{\scriptsize\raggedright
\begin{longtable}{p{0.237\linewidth}p{0.237\linewidth}p{0.237\linewidth}p{0.237\linewidth}}
\toprule
Frame & Field & Before & After \\
\midrule
\endhead
Alert & node id (offset 2, 4 bytes) & vehicle id in clear & pseudonym p, 4 bytes: the leading 4 bytes of HMAC-SHA256(k\_v, "P" and boot (u32, big endian) and ctr (u32, big endian)) \\
Alert & freshness field, sequence, boot counter, event id & in clear & removed from the clear: the base takes boot and counter from the window entry that p matches; the bytes carry encrypted fields (a stream cipher keyed from k\_v with the nonce boot and ctr) \\
Alert & latitude and longitude (offsets 22 to 29) & in clear & encrypted in place; authenticated by the tag \\
Alert & tag (8 bytes) & HMAC-SHA256(K\_vehicle, bytes 0 to 88)[:8] & HMAC-SHA256(K\_vehicle, bytes 0 to 88 and boot and ctr (8 bytes, not sent))[:8] \\
ACK & node id & vehicle id in clear & the acknowledged pseudonym p \\
Flash and cancel & short id (2 bytes) & vehicle short id in clear & pseudonym p, 2 bytes \\
Flash and cancel & freshness low 16 bits, boot byte & in clear & removed from the clear (14 bytes instead of 15 for the frame with the boot byte); the freed 2 bytes carry encrypted flags and a coarse cell \\
Micro frame & pending index & assigned by the base & unchanged \\
\bottomrule
\end{longtable}}

Base: a hash index from pseudonym to (vehicle, boot, counter) over a window of the next m = 16 counters per vehicle; every matching entry is verified with the tag; the window moves to the matched counter; after a reboot an unresolved frame is tested against the next three boot values. Cost: at most 23 microseconds per frame in Python and about 3 in C for 1,000 vehicles (\texttt{ROUND7\_B.md} section 4); the frame after a reboot costs 84 ms in Python.


\subsection*{10. Round 7b notes}

\textbf{10.1 The micro-confirm index is an unauthenticated lookup hint (G0.3).} The leading byte of the 8-byte micro frame is outside the MAC input. Read from the code (\texttt{crashguard.frame\_layout.micro\_lookup}): the base takes \texttt{pending.get(frame[0])}; if no entry has that index it returns False (reject), and otherwise it verifies the tag under that one entry and rejects on a mismatch. There is no fallback to trial matching on an index miss: \texttt{micro\_lookup\_trial} exists only as the rejected comparison of section 7 and is not on the base path. Therefore:

\begin{itemize}
\item changing the index of a genuine micro frame is the same as dropping it: the frame is rejected, the pending entry is not consumed, and the node's retry sends the full frame (section 6);
\item the forgery bound stays 2\textasciicircum{}-40 per frame (one candidate entry); the n x 2\textasciicircum{}-40 bound of the trial path does not apply, because the trial path is not taken. An attacker who rewrites the index spends the frame on one guess against another entry, which has the same 2\textasciicircum{}-40 chance as any forged frame.
\end{itemize}

The ns-3 frame (13 bytes) looks the entry up by device address (row F7), a different abstraction with the same reject-on-miss behaviour. Evidence: the function above and \texttt{tests/test\_frame\_layout.py} (an unknown index, the index of another entry and a stale boot counter all reject).

\textbf{10.2 Quarantine with pseudonyms: status of the injective no-framing lemma (G3).} What the round-7 report called "not proved" is, exactly, the following. In ProVerif (\texttt{formal/proverif/r7\_lemmas\_quarantine.pv}) the non-injective form (a strike against v implies that v sent a bad alert, or v's key was revealed) is \textbf{proved}, and the injective form is \textbf{reported false with a trace}. The trace (\texttt{formal/out/r7/pv\_r7\_lemmas\_quarantine.txt}) has two copies of the base thread, which both read the seen-set (the else branch of \texttt{get seen}) before either inserts: two strikes for one car event. That is a race between the check and the insert, which the model treats as two steps; it is not a replay across time and not a forgery. It is a requirement on the base: \textbf{the seen-set check and insert must be one atomic step} (one thread per frame stream, or a lock). Round 6 and round 7 had no Tamarin model of the quarantine, so the injective form had no second tool.

Fix and rerun (round 7b, \texttt{formal/r7b\_quarantine\_pseudo.spthy}, Tamarin 1.12.0, outputs in \texttt{formal/out/r7b/}, keys \texttt{r7b.g3.*}): with the seen-set as a linear token consumed by the strike (the atomic step), executability, the non-injective form and the \textbf{injective form are verified} (under one second each). Controls: with no seen-set (the race), the injective form is falsified and the non-injective form is verified, which reproduces the ProVerif result; with the strike before the tag check (rule A), both forms are falsified. Limits: this is a separate small model (the window entry stands for the pseudonym resolution; no window advance, no reboot, no interaction with the alert model), and atomicity is an assumption about the base, not something the model derives.

Lemmas that hold with pseudonyms, and the tool: \textbf{ProVerif}: diff-equivalence for two alerts and for unboundedly many (the fixed-id control and the boot-in-clear control fail); authenticity; unique alert ids with the boot counter (the no-write control fails); micro-confirm and micro-cancel binding (the control without T1 fails); quarantine, non-injective no-framing (the rule-A control fails). \textbf{Tamarin}: L4, L4b, L5, L5 with a car reboot, L7 and both executability lemmas (verified, model \texttt{formal/r7\_alert\_pseudo\_fresh\_bootctr.spthy}); L6, the replay after a base reboot, is refuted as in the base design; quarantine injective and non-injective no-framing with an atomic seen-set (round 7b model above). \textbf{Not covered by any proof}: the quarantine injective form without the atomicity assumption (it fails, by the race), key-holding attackers, timing and radio fingerprints, the window advance and the resynchronisation after a reboot. The paper therefore says that the quarantine rule keeps its no-framing property with pseudonyms when the base applies the seen-set atomically; it does not say more.

\textbf{10.3 Does the tag-covers-boot-and-counter fix touch a frozen file? (G3)} No. The fix (found by the C resolver benchmark, log row B1.fix of round 7) changes only the pseudonym variant: \texttt{src/crashguard/r7\_pseudo.py}, \texttt{firmware\_r7/}, \texttt{formal/proverif/r7\_*.pv} and \texttt{formal/r7\_alert\_pseudo\_fresh\_bootctr.spthy}, all new in round 7. A comparison of \texttt{submit-v3} with \texttt{submit-v4} lists no modified file in \texttt{xsim/}, \texttt{firmware\_rd/}, \texttt{firmware\_r6/}, \texttt{firmware\_export/}, \texttt{formal/} (earlier models) or the frozen modules of earlier rounds; \texttt{scripts/r5\_audit.py} passes. In the base design the tag already covers every byte that precedes it, so the boot counter and the sequence number are covered where they are in the frame; the base design is unchanged. The variant differs only in that boot and counter are not sent and are therefore covered as unsent associated data. Abstraction row F21 (new): the tag input of the pseudonym variant is the frame bytes plus boot and counter of the matched window entry (abstraction: in the models the tag input is a term that contains them); a collision of pseudonyms is decided by this tag.



\subsection*{11. Round 7c: atomic frame handling at the base (design requirement)}

**11.1 What the code does (reading, \texttt{scripts/r7c\_atomic\_scan.py}, keys \texttt{r7c.atomic.*}).** The base-side code is a reference implementation, not a server. \texttt{BaseVerifier.handle} (\texttt{src/crashguard/verifier.py}) does the counter check (\texttt{a.counter <= self.last.get(...)}) and the update (\texttt{self.last[...] = ...}, the \texttt{recent} table, the saved state) as consecutive statements of one method, with no call between them that can yield. The quarantine model (\texttt{r5\_l8.process}) is a pure transition function over one state tuple, and the pseudonym resolver (\texttt{r7\_pseudo}) is sequential methods on one object. A scan of these modules and of every C or C++ source of \texttt{firmware\_rd}, \texttt{firmware\_r6}, \texttt{firmware\_r7}, \texttt{firmware\_export} and \texttt{xsim/ns3} (24 files) finds \textbf{no lock, mutex, thread, asynchronous call, atomic or interrupt primitive} (the only parallel code runs independent experiments, each with its own state; two \texttt{volatile} sinks in the footprint anchor keep the compiler from removing code and are not synchronisation). So the answer is: \textbf{neither a lock nor a threaded design; the check and the insert are atomic only because the reference code is called from one sequential loop. Nothing enforces it}, there is no base server in the repository, and the steps resolve, tag check, seen-set check, seen-set insert and strike are not one code path anywhere (the verifier holds the counters, the quarantine lives in the separate exploration model).

\textbf{11.2 Requirement R-ATOMIC (design requirement for any base implementation).} For every received frame the base executes pseudonym resolution, tag verification, the seen-set or counter check, the seen-set insert and the strike count as one critical section per resolved vehicle: either one single-threaded frame loop that every gateway feeds (a frame forwarded by two gateways is the same frame and passes through the same loop), or a lock per resolved vehicle held from the check to the insert and the strike, with the duplicate-frame path under the same lock. The insert must happen before the lock is released and before any later frame of that vehicle is looked at. A base that reads the seen-set on several threads or processes without this has the race of the ProVerif trace: two strikes for one car event.

\textbf{11.3 Why it matters.} The quarantine lemma with pseudonyms (section 10.2) holds in its injective form only with R-ATOMIC (Tamarin, ATOMIC variant: verified; NOSEEN variant, the race: falsified; ProVerif: trace of two base copies). The non-injective form (no vehicle is struck unless it sent a bad alert or its key was revealed) does not depend on it. R-ATOMIC changes no frame field and no byte of the layout; it constrains the base's implementation only. Abstraction row F22 (new): the seen-set check and insert are one transition in the Tamarin model and two steps in the ProVerif model; the Tamarin result is the one that applies to a base that meets R-ATOMIC, and the ProVerif trace applies to one that does not.



\subsection*{12. Round 7d: re-acknowledgement limits at the base (deployment requirement R-REACK)}

\textbf{12.1 The requirement.} A base that acknowledges alert frames must meet R-REACK: (a) it re-acknowledges the same (node, event) no sooner than a \textbf{re-ACK interval} after its last ACK of that pair, and (b) it pays every ACK, new or repeated, from a \textbf{downlink budget}, a token bucket that refills at the duty cycle, and sends no ACK when the bucket is empty (it still accepts a new valid alert; it just cannot acknowledge it). Both limits are required, and they apply to repeats of a valid frame; a frame whose tag fails never reaches them, because it gets no ACK at all (section 11).

\textbf{12.2 Why it is needed.} The tag proves that a frame was made with the vehicle's key, not that it is new. A node without the key that records a valid alert can send it again, and the base, which re-acknowledges a repeat so that a node whose ACK was lost can stop retrying, answers every time. \textbf{With both limits off, a replayed valid frame is re-acknowledged each time it arrives, and the re-ACKs consume the base's downlink budget.} \texttt{BaseVerifier} (\texttt{src/crashguard/verifier.py}) has both limits OFF by default (\texttt{reack\_interval\_s=None}, \texttt{budget=None}), so the default does not meet R-REACK.

**12.3 Values used (stored replay-flood study, \texttt{crashguard.attacks.flood}, keys \texttt{attacks.flood.*}; regression in \texttt{crashguard.r7d\_base}, keys \texttt{r7d.reack.*}).**

{\scriptsize\raggedright
\begin{longtable}{p{0.475\linewidth}p{0.475\linewidth}}
\toprule
Item & Value \\
\midrule
\endhead
Re-ACK interval & the shortest spacing of two tries of the same alert, \texttt{delivery\_sim.retry\_interval\_s(SF)} = alert airtime + 0.1 s processing + ACK airtime + 0.1 s margin: 0.43 s at SF7, 0.61 s at SF8, 0.94 s at SF9. A legitimate retry cannot arrive sooner, so the limit blocks no legitimate retry \\
Downlink budget & token bucket, refill at the duty cycle \texttt{crashguard.lora.DUTY} = 10\%, capacity 360 s; the study starts it at 30 s (the rest used by other traffic) \\
Flood & one captured valid alert replayed 20 times a second for 120 s at SF7, a genuine alert at 100 s \\
\bottomrule
\end{longtable}}

{\scriptsize\raggedright
\begin{longtable}{p{0.317\linewidth}p{0.317\linewidth}p{0.317\linewidth}}
\toprule
Result (SF7, 26-byte ACK) & Limits off & R-REACK on \\
\midrule
\endhead
Replays re-acknowledged & 679 & 266 \\
Downlink airtime of the re-ACKs & 41.9 s & 16.4 s \\
Replays refused for an empty budget & 1,721 & 0 \\
Budget left at the end & 0.05 s & 25.4 s \\
Genuine alert at 100 s & accepted, \textbf{no ACK} & accepted and acknowledged \\
\bottomrule
\end{longtable}}

Reading. The limits stop the starvation of genuine ACKs; they do not stop the waste: with R-REACK on, the replays still cost 16.4 s of airtime in 120 s (55\% of the 30 s budget at the start, before refill), because the interval is only the retry spacing. A larger interval, or a limit on re-ACKs per pair, would cut it further; that was not evaluated (NOT\_RUN). The values are those of a simulation of the reference code, not of a base server.

\textbf{12.4 Reference configuration.} \texttt{crashguard.r7d\_base.reference\_base(key, sf)} builds a \texttt{BaseVerifier} with both limits on and the values above. It is a new module; \texttt{verifier.py} and every frozen file are unchanged, and the default of \texttt{BaseVerifier} stays off. The earlier flood study already built its limited base the same way (\texttt{attacks.py}). \texttt{crashguard.r7d\_base.main} reruns the flood with the limits off and on and compares it with the stored \texttt{attacks.flood.*} keys (both match). Abstraction row F23 (new): the ns-3 gateway and the Tamarin and ProVerif models have no re-ACK limit; the requirement is outside the models.


